\section{Complete metrics and the auxiliary end factor}
\label{sec:rpi:lee-completion}

Lee's quantitative positive-mass estimate turns small mass into control
of a harmonic end factor. To apply it to $J_t$, we construct complete
smooth metrics whose masses tend to zero and whose end factors approach
$J_t$. The construction keeps a fixed coordinate tail while allowing
the compact interior to vary. Throughout, $n\ge8$ and $d=n-2$.

\begin{proposition}[Complete approximants and the auxiliary factor]
\label{prop:rpi:complete-end-factors}
For the normalized flow of Proposition~\ref{prop:rpi:mass-evolution}
there are fixed radii $R_c<R_1$
and complete smooth connected boundaryless one-ended metrics with nonnegative scalar
curvature, harmonic end factors $\mathcal Z_t$ on $\{r>R_1\}$,
and masses $\mathcal M_t$, such that
\[
 0\le\mathcal M_t\longrightarrow0,\qquad
 \sup_{r\ge2R_1}r^d|\mathcal Z_t-J_t|\longrightarrow0.
\]
The radius $R_c$ encloses every normalized frontier and compact core.
For every fixed $s>R_c$,
\begin{equation}\label{eq:rpi:auxiliary-weighted-limit}
 \sup_{r\ge s}r^d|J_t-1|\longrightarrow0.
\end{equation}
\end{proposition}

We first isolate the quantitative input, so that the construction can
be read separately from the verification of its positive-mass hypotheses.

\begin{lemma}[Quantitative positive mass on one harmonic end]
\label{lem:rpi:quantitative-one-end}
For each $n\ge8$ there is a fixed constant $\ell_n>1$ such that,
for every $\eta>0$, there is $\delta(n,\eta)>0$ with the following
property. Let $(N,h)$ be a smooth complete connected boundaryless
$n$-manifold with $R_h\ge0$. Suppose that outside a compact set it
consists of a single coordinate end $\{r>R\}$, where $R>0$, and that
on this end
\[
 h=Z^{4/(n-2)}\delta,\qquad Z>0,\qquad
 \Delta_\delta Z=0,\qquad Z\longrightarrow1.
\]
Writing $d=n-2$ and $m=m_{\rm ADM}(h)$, we have
\[
 m<\delta(n,\eta)R^d
 \quad\Longrightarrow\quad
 \sup_{r>\ell_nR}|Z-1|<\eta.
\]
\end{lemma}

\begin{proof}
Rescale the metric by $R^{-2}$ and the coordinates by $R^{-1}$.
The harmonic coordinate radius becomes one, and the mass becomes
$m/R^d$. We apply the argument of
\cite[Theorem~1.3]{Lee2007NearEquality} at this scale. Its formal
hypothesis concerns every complete asymptotically flat metric on the
underlying manifold; we verify that the positive-mass statements used
in its proof follow here from Lemma~\ref{lem:rpi:harmonic-pmt}.
The one-end assumption also identifies the mass flux in Lee's
Lemma~2.1 with the flux in the divergence integral over the entire
manifold.

The scalar-flat reduction of
\cite[Lemma~2.5]{Lee2007NearEquality} gives a metric
$k=f^{4/d}h$ with $f>0$ and $f\to1$. The factor is bounded above
and below by positive constants, so $k$ remains smooth and complete.
Its end factor $Zf$ is Euclidean harmonic. For this scalar-flat
reference metric, Lee uses the variations
\[
 k_s=k+s\zeta\operatorname{Ric}(k),\qquad
 \widehat k_s=u_s^{4/d}k_s,
\]
where the cutoff $\zeta$ is compactly supported and
$\widehat k_s$ is scalar-flat. Outside the support of $\zeta$, the
product of the old end factor and $u_s$ is Euclidean harmonic.
The conformal scalar-curvature identity gives
\[
 c_n=\frac{n-2}{4(n-1)},\qquad
 L_s=\Delta_{k_s}-c_nR_{k_s},\qquad
 L_s(u_s-1)=c_nR_{k_s}.
\]
Both scalar-curvature coefficients in the operator definition and
inhomogeneous equation preceding Lee's Lemma~2.4 must have this value,
as in the earlier conformal equation in that paper. Lemma~2.4 estimates
$\Delta_k$, and the ensuing perturbation argument applies with the
corrected fixed dimensional coefficient. Its exterior estimates give
$u_s-1=O(|s|)$ uniformly. On the compact region away from the
perturbation, $u_s-1$ is harmonic, so the maximum principle extends
this bound throughout the manifold. For sufficiently small $|s|$,
$1/2\le u_s\le3/2$. The perturbed and conformally corrected metrics
are therefore smooth and complete.

The scalar-flat argument uses nonnegativity of the varied mass at a
small negative parameter
\cite[proof following Lemma~2.2]{Lee2007NearEquality}. This follows
from Lemma~\ref{lem:rpi:harmonic-pmt}, since every varied metric has
a harmonic end outside a compact set. The same lemma gives
$m_{\rm ADM}(k)\ge0$ after the initial scalar-flat reduction.
Lee's Lemma~2.5 gives $m_{\rm ADM}(k)\le m_{\rm ADM}(h)$ and
controls the end-factor difference by the mass difference.
Thus all the required
positive-mass inputs hold, without a spin or orientability assumption.
The constants in Lee's estimates depend only on dimension and the
target error at the normalized radius, not on the compact interior.
Restoring the scale $R$ gives the asserted implication.
\end{proof}

\begin{proof}[Proof of Proposition~\ref{prop:rpi:complete-end-factors}]
We use the following elementary harmonic estimate.  If $q\ge0$ is
Euclidean harmonic on $\{r>R\}$, tends to zero, and has expansion
$q=b r^{-d}+O(r^{-d-1})$, then
\begin{equation}\label{eq:rpi:positive-harmonic-monopole}
 0\le q(x)\le C(n)b r^{-d}\qquad(r\ge2R).
\end{equation}
Indeed, its spherical average is exactly $b r^{-d}$, and the Harnack
inequality on scaled annuli bounds its maximum on a sphere by a
dimensional constant times this average.  In particular, the constant
does not depend on the geometry inside $B_R$.

\paragraph{A smooth reduction from two ends to one.}
Let $(N,k)$ be a smooth complete connected boundaryless manifold with
exactly two harmonically flat ends and compact end complement, and
suppose $R_k\ge0$. Designate one end for retention. There is a
harmonic function $\chi$ tending to one on that end and zero on the
other, with $0<\chi<1$.  To construct it, choose a smooth function
constant with these values outside a compact set and minimize its
Dirichlet energy after addition of the homogeneous energy space.
The global Sobolev inequality gives the solution; truncation and the
strong maximum principle give its bounds.  Exterior elliptic estimates
give the stated end limits and expansions.

Write $k=Z_i^{4/d}\delta$ on either end.  Then $Z_i\chi$ is Euclidean
harmonic.  On the end to be filled,
\[
 Z_2\chi=b_2r^{-d}+O_j(r^{-d-1}),\qquad b_2>0.
\]
Positivity of $b_2$ follows by comparison with a positive multiple of
$r^{-d}$ from any fixed coordinate sphere.  Under inversion
$x=y/|y|^2$, the factor
\[
 |y|^{-d}(Z_2\chi)(y/|y|^2)
\]
extends harmonically across $y=0$ with value $b_2$.  Thus
$\chi^{4/d}k$ extends to a smooth metric after adding one point.
The filled manifold is complete, has one end, and has nonnegative
scalar curvature, since
$R_{\chi^{4/d}k}=\chi^{-4/d}R_k$ away from the added point and is zero
near that point.

At the retained end write
$\chi=1-A r^{-d}+O_j(r^{-d-1})$.  If $M=m_{\rm ADM}(k)$, the filled
metric has mass $M-2A$.  The function $Z_1(1-\chi)$ is nonnegative
Euclidean harmonic with monopole $A$.  Consequently
\begin{equation}\label{eq:rpi:one-end-reduction}
 0\le M-2A\le M,\qquad
 |Z_1-Z_1\chi|\le C(n)M r^{-d}\quad(r\ge2R),
\end{equation}
where the first inequality uses
Lemma~\ref{lem:rpi:harmonic-pmt} on the smooth filled manifold.
The estimate depends only on the retained-end coordinate radius and
mass; the second end may vary with the compact interior.

\paragraph{A common coordinate tail.}
Use the normalized metric $G_t$ and total harmonic factor $F_t$
from Section~\ref{sec:rpi:flow-mass-capacity}. In this subsection,
$\Sigma_t,w_t$ denote its frontier and capacitary potential in these
coordinates; the latter is the pullback previously written
$w_t\circ\Phi_t$. A coordinate ball includes the compact core when
the whole exterior is under discussion. The mass, capacity, and deficit
remain $m(t)$, $c_t$, and $\mu(t)=m(t)-c_t$.
The preceding radius bound and harmonic estimates give
$\Sigma_t\subset B_{R_c}$, the factors $F_t$ are uniformly bounded on
a fixed sphere outside $B_{R_c}$, and
\[
 0\le c_t\le m(t)\le m(0),\qquad \mu(t)\longrightarrow0.
\]
Indeed, take $T=t-1$ in the radius estimate: $U_t\le4e$ on the
normalized tail $r\ge R_{\max}$ for $t\ge1$. The initial bounded time
interval is covered by the finite-time flow bounds. Multiplication by
the uniformly bounded initial harmonic factor gives the asserted
bound for $F_t$. Increasing $R_c$ once includes that sphere. Let $w_t$ be the
capacitary potential equal to one at $\Sigma_t$ and zero at infinity.

For $0<\lambda<1$ put $p=1+\lambda w_t$ and
$h_{t,\lambda}=p^{4/d}G_t$.  The detachment and smoothing construction
above provides a smooth enclosing boundary $\Gamma$ with
$H_\nu^\Gamma<0$.  It can be chosen in a ball $B_{R_1}$ whose radius
is independent of $t$ and $\lambda$.  Here is the confinement
argument.  On the common coordinate end the positive harmonic factor
$F_tp$ has a uniform bound on $S_{R_c}$ and tends to one.  The maximum
principle and scaled gradient estimates give
\[
 |F_tp-1|+r|D(F_tp)|\le C(R_c/r)^d\quad(r\ge2R_c).
\]
Thus all coordinate spheres beyond a fixed $R_*$ have positive
outward mean curvature in $h_{t,\lambda}$. Write $K_t$ for the
original conductor and $L$ for its detached minimizing enlargement.
The functional minimized by $L$ is
\[
 \mathcal F(L)=P_{h_{t,\lambda}}(L)+
       \theta\int_{L\setminus K_t}h_0\,dV_{h_{t,\lambda}},
 \qquad h_0\ge0 .
\]
Let $X$ be the outward unit normal to the coordinate spheres.
For almost every $s>R_*$, the finite-perimeter divergence formula
on $L\cap\{r>s\}$ gives
\[
 P(L\cap\{r<s\})
 \le P(L)-\int_{L\cap\{r>s\}}\operatorname{div}X\,dV.
\]
Both the perimeter and the nonnegative volume penalty decrease under
clipping, strictly if the removed tail has positive volume.  Minimality
therefore confines $L$ to $B_{R_*}$.  The subsequent smooth boundaries
lie in shrinking collars, so they lie in $B_{R_1}$ for one
fixed $R_1>R_*$.  Fix such a smoothing index before doubling.
Thus the coordinate tail is fixed even when the separation distance
and smoothing parameters vary.

\paragraph{The capacity comparison and the completed end.}
Suppress $t$ temporarily and put $h=G_t$, $F=F_t$, $w=w_t$, $c=c_t$,
$m=m(t)$, $\mu=\mu(t)$.  The potential and capacity of the original
conductor for $h_\lambda=p^{4/d}h$ are exactly
\begin{equation}\label{eq:rpi:perturbed-capacity-exact}
 w_{\Sigma,\lambda}=\frac{(1+\lambda)w}{p},\qquad
 c_{\Sigma,\lambda}=(1+\lambda)c,\qquad
 m_\lambda=m+2\lambda c.
\end{equation}
These identities follow from the harmonic conformal transformation
formula and the end expansions.  If $w_\Gamma$ is the potential of
$\Gamma$ for $h_\lambda$ and $c_\Gamma$ its capacity, enclosure
monotonicity and Lemma~\ref{lem:rpi:mass-capacity} give
\[
 0\le D_\Gamma:=c_\Gamma-c_{\Sigma,\lambda}
 \le m_\lambda-c_{\Sigma,\lambda}=\mu+\lambda c.
\]
On the common tail the function
$Fp(w_\Gamma-w_{\Sigma,\lambda})$ is nonnegative Euclidean harmonic
with monopole $D_\Gamma$.  Its sign follows from the maximum
principle on the exterior of $\Gamma$.

For $0<a<1$ the positive two-ended doubling construction in
Lemma~\ref{lem:rpi:mass-capacity} has retained end factor
\[
 B=Fp\left(1-\frac{1-a}{2}w_\Gamma\right)
\]
and retained end mass
\[
 m_B=m_\lambda-(1-a)c_\Gamma
 \le \mu+\lambda c+a m_\lambda .
\]
The Miao smoothing and its conformal correction produce smooth
complete two-ended metrics with $R\ge0$ and retained factor $Bz_\delta$
on $\{r>R_1\}$.  In the notation of that proof,
$z_\delta\ge1$, and
\[
 e_\delta:=B(z_\delta-1)\ge0
\]
is Euclidean harmonic on the common tail with monopole
$A_\delta\to0$.  Their retained masses are
$M_\delta=m_B+2A_\delta\ge0$.

The auxiliary Euclidean end factor already defined above is
$J=F(1-w/2)$, with $t$ suppressed. Direct subtraction gives
\[
 B-J=
 \frac{\lambda}{2}Fw
 -\frac12Fp(w_\Gamma-w_{\Sigma,\lambda})
 +\frac a2Fp w_\Gamma .
\]
Each of the three harmonic functions appearing on the right is
nonnegative before its displayed sign, with monopoles respectively
$c$, $D_\Gamma$, and $c_\Gamma$.
Equation~\eqref{eq:rpi:positive-harmonic-monopole} therefore yields
\begin{equation}\label{eq:rpi:completed-factor-comparison}
 |Bz_\delta-J|
 \le C(n)\bigl(\mu+\lambda c+a m_\lambda+A_\delta\bigr)r^{-d},
 \qquad r\ge2R_1 .
\end{equation}
The capacity difference controls the end-factor error independently
of the distance between $\Gamma$ and $\Sigma_t$.

For each $t$, choose $\lambda_t=a_t=(t+2)^{-1}$; choose the
detachment and boundary smoothing parameters next; finally choose
$\delta_t$ so that $A_{\delta_t}\le(t+2)^{-1}$.
The Sobolev constants used in the seam correction may depend on
these fixed choices.  This order suffices because Lee's constants
do not depend on the compact interior.  The preceding bounds imply
\[
 0\le M_{\delta_t}\longrightarrow0,\qquad
 \sup_{r\ge2R_1}r^d|B_tz_{\delta_t}-J_t|\longrightarrow0.
\]
Apply the one-end reduction
\eqref{eq:rpi:one-end-reduction} to each smooth completed metric.
Its underlying double is connected because the original exterior is
connected and $\Gamma$ is nonempty; it is boundaryless with exactly
two ends and compact end complement.
Denote the resulting masses and harmonic end factors by
$\mathcal M_t$ and $\mathcal Z_t$. These smooth connected one-ended
metrics have their harmonic coordinate end on the same $\{r>R_1\}$,
and the first two assertions of the proposition follow from
Equation~\eqref{eq:rpi:one-end-reduction}.
Lemma~\ref{lem:rpi:quantitative-one-end}, followed by the
end-factor comparison, now gives
\[
 J_t\longrightarrow1
 \quad\hbox{uniformly on }\{r>R_2\}
\]
for any one fixed $R_2>\max(2,\ell_n)R_1$.
The exterior maximum principle also gives weighted convergence on
a slightly larger fixed tail.

\paragraph{The normalization and the inner tail.}
The initial harmonic factor in the normalized coordinates is
$h_t=H\circ\Phi_t$. Thus
\[
 \bar G_t=h_t^{4/d}\delta,\qquad
 F_t=h_tU_t,\qquad J_t=h_tW_t,\qquad
 W_t=\tfrac12(U_t-V_t).
\]
The initial end expansion gives
$h_t-1=O_j(e^{-2t}r^{-d})$ on any fixed exterior tail.
Consequently the preceding conclusion implies $W_t\to1$ there.

It remains to extend the weighted convergence to each fixed radius
$s>R_c$. The functions $J_t$ are positive Euclidean harmonic
on $\{r>R_c\}$, with spherical averages
$1+\tfrac12\mu(t) r^{-d}$.  Their bounded masses therefore give
uniform Harnack and elliptic bounds on each compact subannulus.
Every subsequential harmonic limit equals one on the distant tail,
and hence everywhere on $\{r>R_c\}$ by unique continuation.
Thus $J_t\to1$ uniformly on $S_s$.  Applying the maximum principle
to $J_t-1$ outside $S_s$ yields
\[
 \sup_{r\ge s}r^d|J_t-1|
 \le s^d\sup_{S_s}|J_t-1|\longrightarrow0.
\]
The same conclusion holds for $W_t$, since $h_t\to1$ with the
stated decay. This proves Equation~\eqref{eq:rpi:auxiliary-weighted-limit}.
\end{proof}
